\documentclass[../../main.tex]{subfiles}
\begin{document}
\begin{oframed}
  \subfile{problem}
\end{oframed}

{\bf [解答]}

$A_n = \sum_{i=1}^n a_i$ とおくと，題意より$i=1,2,\cdots,n$ に対し
\begin{align}
  a_i \le \dfrac{A_n - a_i}{n-1} \label{eq:1}
\end{align}
だから，$i=1,2,\cdots, n$ について足し合わせて$n-1 > 0$ から
\begin{align}
  (n-1)A_n \le (n-1)A_n
\end{align}
となる．よってこの等号が成立する．
つまり，\cref{eq:1}の等号が全て成立するから整理して
\begin{align}
  a_i = \dfrac{A_n}{n}
\end{align}
であるから
\begin{align}
  a_1 = a_2 = \cdots = a_n
\end{align}
の時題意がみたされる．

\end{document}
